\chapter{Conclusions and Future Work}
\label{ch:conclusions}

\section{Achievement of the Research Objectives}
The first objective was to obtain a model suitable for controller design. A
first-principles nonlinear model was derived, specialized using measured
geometry and mass properties, and reduced to the upright relation
$\ddot{\alpha}=100.8\alpha-1.952\ddot{\theta}$. The complete derivation is
retained in Appendix~\ref{app:full_derivation}.

The second objective was to construct an integrated experimental stack. The
completed platform measures both angles, estimates angular rates, applies
wrap-safe state processing, integrates acceleration commands into step rate,
enforces travel and nonlinear-validity limits, records trial data, and
exports auditable figures and tables.

The third objective was to implement three controllers on a common platform.
Linear full-state feedback, hybrid sliding-mode control, and four-state
sliding-mode control were all implemented through the same sensing and
actuation path. The work also demonstrated commanded base-reference motion
with the linear controller.

The fourth objective was experimental comparison. Seven fixed trials were
analysed across hold, nudge, and tap conditions. This objective was achieved
as a qualitative engineering comparison, not as a statistical evaluation.

\section{Answers to the Research Questions}
The acceleration-input representation provides a useful bridge between the
Furuta model and the firmware command structure, provided it is described as
an abstract command. The real stepper can still stall or miss motion, so
measured arm feedback and actuator diagnostics remain necessary.

For the pinned hold trials, the linear controller produced the lowest RMS
pendulum error, the hybrid SMC produced the smallest peak pendulum and arm
excursions, and SMC4 required the greatest RMS command effort. In the pinned
tap trials, only the linear controller completed the run without falling.
These observations describe the selected runs only and do not establish a
general ranking of linear and sliding-mode control.

The practical mechanisms required for usable hardware control were filtered
derivatives, wrap-safe angle differences, sensor-jump rejection, calibration
and sign checks, low-speed deadband, stop-before-reverse behavior, engagement
ramping, soft and hard travel limits, nonlinear validity guards, and
stall/missed-motion diagnostics.

\section{Main Findings}
\begin{itemize}
  \item The linear controller was the most consistent overall controller in
  the pinned dataset and remained upright in the representative tap run.
  \item The hybrid SMC reduced peak excursions in the representative hold run,
  but its representative tap recovery ended in a fall.
  \item SMC4 balanced in the representative hold run but was the most
  sensitive to actuator limits, denominator protection, and command shaping.
  \item The nudge run demonstrated balance with commanded arm motion, but none
  of the three steps met the fixed $\pm\SI{1}{\degree}$ settling criterion
  within the available windows.
\end{itemize}

\section{Limitations}
Each condition is represented by one pinned trial. Manual tap magnitude and
onset were not measured, the telemetry rate was lower than the control rate,
and the nominal model was not validated through a dedicated parameter
identification campaign. The work is restricted to upright stabilization and
does not address swing-up. Some component procurement details were also not
retained, which limits exact hardware duplication despite the documented
functional architecture.

\section{Future Work}
\begin{itemize}
  \item repeat every controller--condition combination enough times to report
  distributions, confidence intervals, and failure rates;
  \item add a repeatable disturbance device or trigger so recovery time and
  disturbance magnitude are measurable;
  \item perform dedicated system identification and compare predicted and
  measured trajectories before redesigning controller gains;
  \item log selected controller and actuator signals closer to the
  \SI{200}{Hz} control rate;
  \item investigate actuator-aware nonlinear control, observers, and
  anti-windup or command-allocation methods for bounded stepper motion;
  \item add and evaluate a separate swing-up controller if full operating-range
  control becomes a project requirement.
\end{itemize}
